\documentclass[sigconf,nonacm]{acmart}

\usepackage{listings}
\usepackage{booktabs}
\usepackage{microtype}
\usepackage{xcolor}
\usepackage{url}

\lstset{
  basicstyle=\ttfamily\footnotesize,
  breaklines=true,
  frame=single,
  columns=fullflexible,
  backgroundcolor=\color{gray!8},
  rulecolor=\color{gray!40}
}

%% -----------------------------------------------------------
%% Preamble metadata (acmart reads these before \begin{document})
%% -----------------------------------------------------------
\title{WiFi-Based Automated Classroom Attendance
       Using a Commodity Android Smartphone}

\author{Vivek Barhate}
\affiliation{%
  \institution{George Mason University}
  \department{Department of Computer Science}
  \country{USA}}

\author{Preshita Bhortake}
\affiliation{%
  \institution{George Mason University}
  \department{Department of Computer Science}
  \country{USA}}

\begin{abstract}
Manual attendance in large classrooms is slow, error-prone, and
vulnerable to proxy marking.  This paper presents a self-contained,
privacy-preserving attendance system built from a single commodity
Android smartphone acting as a mobile WiFi hotspot.  Students connect
to the hotspot and check in through a captive-portal web page ---
no dedicated app is required on the student device.  An ICMP
ping-sweep engine running on the professor's phone continuously tracks
device associations and measures presence duration; JavaScript
heartbeats from the student browser provide an application-layer
presence signal.  A SHA-256 token derived from the client IP and a
per-session random salt replaces every raw IP and MAC address in the
database, providing strong pseudonymisation.  Two server-side guards
prevent proxy check-in: each device may check in only once per
session, and each roll number may appear only once per session.  A
FastAPI/SQLite backend provides session management, live occupancy,
and CSV export.  End-to-end tests confirm 8/8 router-app flows and
12/12 backend flows; a 35-student concurrent load test shows median
heartbeat latency below 15~ms.
\end{abstract}

\keywords{WiFi attendance, captive portal, Android hotspot, ping sweep,
  privacy, NanoHTTPD, FastAPI}

%% -----------------------------------------------------------
\begin{document}
\maketitle
%% -----------------------------------------------------------

%% -----------------------------------------------------------
\section{Introduction}
%% -----------------------------------------------------------

Taking attendance in large university classrooms is a persistent
administrative burden.  Sign-in sheets are trivially forged through
proxy marking.  QR-code systems require individual scanning, which
disrupts the start of class.  Dedicated RFID or Bluetooth hardware
requires procurement budgets that many institutions cannot sustain.

WiFi offers a compelling alternative: every student already carries a
smartphone, and connecting to the classroom network is a natural,
low-friction indicator of physical presence.  Prior work has explored
802.11 probe-request sniffing~\cite{demir2014} and router association
tables~\cite{hwang2017}; however, these approaches require
infrastructure-mode access points or passive monitoring hardware that
collects raw MAC addresses, raising privacy concerns.

This project builds a fully self-contained attendance system with the
following properties.

\begin{itemize}
  \item Runs on \emph{one} commodity Android smartphone --- the
        professor's phone acts as both the WiFi hotspot and the
        attendance server.
  \item No dedicated app is needed on the student's device; any
        smartphone browser suffices.
  \item No internet connection is required during class.
  \item No IP or MAC address is ever stored; all tracking uses
        short-lived, session-scoped anonymous tokens.
  \item Two independent server-side guards prevent proxy check-in
        without any additional hardware.
\end{itemize}

%% -----------------------------------------------------------
\section{System Architecture}
%% -----------------------------------------------------------

The system has three logical tiers: the \emph{Router App} on the
professor's phone, the \emph{FastAPI backend} on a laptop on the same
LAN, and the student's standard \emph{web browser}.

\subsection{Data Flow Overview}

The sequence below summarises an end-to-end session.

\begin{enumerate}
  \item The professor starts a session in the Router App, which spins
        up an Android \texttt{LocalOnly\-Hotspot} and the embedded
        NanoHTTPD server on port 8080.
  \item The student connects to the hotspot.  On many devices the OS
        opens the captive-portal page automatically; otherwise the
        student navigates to \texttt{http://<hotspot-ip>:8080}
        manually.
  \item The student submits their roll number.  The server validates it
        against the roster, runs the two anti-proxy guards, computes a
        session-scoped token, writes a \texttt{check\_in} row to the
        local Room database, and redirects to a success page.
  \item The success page runs a JavaScript interval that POSTs
        \{\texttt{token}, \texttt{session\_id}\} to
        \texttt{/heartbeat} every 30~s, keeping a \texttt{presence\_log}
        row updated as long as the page remains open.
  \item In parallel, a ping-sweep coroutine scans the hotspot subnet
        every 5~s, detects associated IPs, and upserts
        \texttt{presence\_log} rows for every responding device.
  \item At session end, \texttt{AttendanceEngine} joins
        \texttt{check\_in} and \texttt{presence\_log} to compute the
        attendance status for every student in the roster.
\end{enumerate}

\subsection{Router App (Professor's Phone)}

The Router App is a Kotlin/Android application.  It uses
\texttt{WifiManager.LocalOnlyHotspotCallback} to create a local-only
WiFi hotspot.  Android assigns a random SSID
(e.g., \texttt{AndroidShare\_XXXX}) and subnet dynamically; the app
reads the actual IP prefix from the \texttt{LocalOnlyHotspotConfiguration}
object at runtime.

\paragraph{Embedded HTTP server.}
NanoHTTPD~\cite{nanohttpd} is bundled as an in-process HTTP server on
port 8080.  Table~\ref{tab:router-endpoints} lists the endpoints.

\begin{table}[h]
\caption{Router App HTTP endpoints (port 8080)}
\label{tab:router-endpoints}
\begin{tabular}{lll}
\toprule
Method & Path & Purpose \\
\midrule
GET  & /                  & Check-in form \\
POST & /checkin           & Validate roll number, create check-in \\
GET  & /success           & Confirmation page with JS heartbeat \\
POST & /heartbeat         & JS keep-alive from browser \\
GET  & /session           & Active session JSON (diagnostic) \\
GET  & /ping              & Connectivity probe \\
GET  & /generate\_204     & OS captive-portal probe (204) \\
GET  & /hotspot-detect.html & iOS captive-portal probe (204) \\
\bottomrule
\end{tabular}
\end{table}

\paragraph{Ping-sweep presence engine.}
Android 10+ blocks \texttt{/proc/net/arp} for third-party apps,
preventing direct ARP-table reads.  \texttt{HotspotManager} runs ICMP
ping sweeps across the full hotspot subnet on a background coroutine.
The sweep interval is 5~s during an active session and relaxes to 60~s
when idle (\emph{adaptive sampling}).  Each ping response yields an
(IP, RTT) pair.  The IP is anonymised immediately:
%
\begin{equation}
  \text{token} = \text{SHA-256}(\text{clientIP} \;\|\; \text{sessionSalt})
\end{equation}
%
where \texttt{sessionSalt} is a 128-bit random value generated at
session start.  The raw IP is discarded; only the token is stored.

\paragraph{Fault tolerance.}
A per-device consecutive-miss counter prevents spurious ``departed''
events.  A device is pruned from the active set only after three
consecutive missed pings ($\approx$15~s at the 5~s interval),
tolerating short-lived radio interference or DHCP renewal gaps.

\paragraph{Signal-quality proxy.}
Android does not expose RSSI for associated hotspot clients.  Ping RTT
serves as a proxy: higher RTT correlates with poor link quality or
greater distance from the hotspot.  The app maintains a rolling window
of the last 10 RTT samples per device and stores the median.  This
value appears in the exported CSV as an informational column.

\paragraph{Attendance thresholds.}
\texttt{AttendanceEngine.kt} classifies each roster student using the
effective time the student's device was present (capped at the session
duration and clamped to start at check-in time):

\begin{itemize}
  \item \textbf{PRESENT} --- $\geq 75\%$ of session duration.
  \item \textbf{PARTIAL} --- $50\text{--}74\%$.
  \item \textbf{ABSENT} --- $< 50\%$ or no check-in at all.
\end{itemize}

\subsection{FastAPI Backend}

A Python REST API (FastAPI~\cite{fastapi} + SQLAlchemy + SQLite)
running on the professor's laptop exposes the endpoints in
Table~\ref{tab:backend-endpoints}.  The backend uses the same
token-centric data model and can be scaled independently.  A UDP
bridge (\texttt{udp\_bridge.py}) accepts heartbeats as lightweight UDP
datagrams on port 9876, reducing HTTP overhead on congested networks.

\begin{table}[h]
\caption{FastAPI backend endpoints (port 8000)}
\label{tab:backend-endpoints}
\begin{tabular}{lll}
\toprule
Method & Path & Description \\
\midrule
POST  & /sessions              & Create session \\
PATCH & /sessions/\{id\}/end   & Close session \\
POST  & /enroll                & Bind token to student identity \\
POST  & /heartbeat             & Record device presence \\
GET   & /occupancy/\{id\}      & Live count (90~s window) \\
GET   & /report/\{id\}         & JSON attendance report \\
GET   & /report/\{id\}/csv     & CSV download \\
GET   & /dashboard             & Live HTML dashboard \\
\bottomrule
\end{tabular}
\end{table}

%% -----------------------------------------------------------
\section{Privacy Design}
%% -----------------------------------------------------------

Privacy is a first-class requirement.  Table~\ref{tab:privacy}
summarises what the system stores and what it deliberately omits.

\begin{table}[h]
\caption{Privacy properties}
\label{tab:privacy}
\begin{tabular}{ll}
\toprule
\textbf{Not stored} & \textbf{Stored} \\
\midrule
IP addresses        & SHA-256 token (session-scoped) \\
MAC addresses       & Roll number / name (self-provided) \\
Location data       & Presence timestamps (first / last seen) \\
Device fingerprints & Session metadata (course, start / end) \\
Browser / OS info   & Median ping RTT (signal quality) \\
\bottomrule
\end{tabular}
\end{table}

The token is computed \emph{server-side} from the client IP and a
random session salt at check-in time.  Because the salt changes each
session, the same physical device produces a different, unlinkable
token in each session.  Identity (roll number, name) is joined to the
token only in the \texttt{check\_in} table; all subsequent presence
records reference the token alone.  The final report joins these tables
in application memory at query time --- the joined result is never
persisted, limiting the scope of a potential database breach.

%% -----------------------------------------------------------
\section{Implementation}
%% -----------------------------------------------------------

\subsection{Android Technology Stack}

\begin{itemize}
  \item \textbf{Room} (SQLite ORM) --- five tables: \texttt{session},
        \texttt{roster}, \texttt{check\_in}, \texttt{presence\_log},
        \texttt{enrollment}.
  \item \textbf{NanoHTTPD 2.3.1} --- embedded HTTP server.
  \item \textbf{Kotlin Coroutines} --- non-blocking ping sweeps
        (\texttt{Dispatchers.IO}) and all DAO calls.
  \item \textbf{Android Foreground Service} --- keeps the hotspot and
        HTTP server alive when the screen is off.
  \item \textbf{Gson} --- JSON parsing in the heartbeat handler,
        wrapped with null-safe Kotlin accessors (see
        Section~\ref{sec:gson-bug}).
\end{itemize}

\subsection{Roster Management}

Before class, the professor imports a CSV roster via
\texttt{Roster\-Activity}:

\begin{lstlisting}
rollNo,name
CS001,Aarav Patel
CS002,Priya Sharma
\end{lstlisting}

A built-in ``Load Sample'' action populates 35 demo students
(CS001--CS035) for testing without a real file.  The roster is stored
in Room and validated at check-in time; a roll number absent from the
roster is rejected with an error shown on the form.

\subsection{Anti-Proxy Guards}
\label{sec:proxy}

The original system had no check-in de-duplication, allowing one device
to enter multiple roll numbers or two devices to check in the same
student.  Two guards were added to \texttt{handleCheckIn} in
\texttt{CheckInHttpServer.kt}.

\paragraph{Guard 1 --- one device per session.}
After computing the token for the incoming request, the server queries:

\begin{lstlisting}[language=SQL]
SELECT * FROM check_in
WHERE sessionId = ? AND token = ? LIMIT 1
\end{lstlisting}

If a row is found, the form is redisplayed with the message
\emph{``This device has already been used to check in for this
session.''} This prevents a student from submitting multiple roll
numbers from the same phone.

\paragraph{Guard 2 --- one roll number per session.}
The server then queries:

\begin{lstlisting}[language=SQL]
SELECT * FROM check_in
WHERE sessionId = ? AND rollNo = ? LIMIT 1
\end{lstlisting}

If found, the form returns \emph{``Roll number X has already been
checked in for this session.''} This prevents two devices from
checking in the same student.

The \texttt{CheckInDao} insert conflict strategy was simultaneously
changed from \texttt{REPLACE} to \texttt{IGNORE} as a database-level
backstop, ensuring a race condition between two concurrent submissions
cannot silently overwrite an existing record.

\begin{table}[h]
\caption{Anti-proxy guard outcomes}
\label{tab:guards}
\begin{tabular}{ll}
\toprule
Scenario & Result \\
\midrule
First check-in: new device + new roll number   & Allowed \\
Same device, different roll number             & Rejected (Guard 1) \\
Different device, same roll number             & Rejected (Guard 2) \\
Different device, different roll number        & Allowed \\
\bottomrule
\end{tabular}
\end{table}

\subsection{Gson Null-Safety Bug and Fix}
\label{sec:gson-bug}

During testing, \texttt{POST /heartbeat} with an empty JSON body
\texttt{\{\}} returned HTTP 500 instead of 400.  Root cause: Gson uses
JVM reflection to construct Kotlin data classes, bypassing Kotlin's
default parameter values and non-nullable type guarantees.  Fields
absent from the JSON payload are set to \texttt{null} at the JVM level
even when declared \texttt{String} (non-nullable) in Kotlin.  A
subsequent \texttt{payload.token.trim()} threw a
\texttt{NullPointerException} outside the inner try-catch, which the
outer \texttt{serve()} handler converted to a 500 HTML response.  The
fix applies null-safe access throughout:

\begin{lstlisting}
// Before (throws NPE when Gson returns null for missing field)
val token = payload.token.trim()

// After
val token = payload?.token?.trim() ?: ""
\end{lstlisting}

%% -----------------------------------------------------------
\section{Evaluation}
%% -----------------------------------------------------------

\subsection{Functional Test Suite}

\texttt{test\_captive\_portal.py} exercises both servers against a live
Android device.  Table~\ref{tab:tests} shows the final results after
all fixes were applied.

\begin{table}[h]
\caption{Functional test results (live device, hotspot IP
  10.116.213.247)}
\label{tab:tests}
\begin{tabular}{llc}
\toprule
Server & Test & Result \\
\midrule
Backend & GET /health returns ok                          & Pass \\
Backend & POST /sessions creates a session                & Pass \\
Backend & POST /enroll enrolls a student                  & Pass \\
Backend & POST /enroll is idempotent (re-enroll updates)  & Pass \\
Backend & POST /heartbeat records presence                & Pass \\
Backend & POST /heartbeat rejects unknown session (404)   & Pass \\
Backend & GET /occupancy returns count $\geq$ 1           & Pass \\
Backend & GET /report returns student in report           & Pass \\
Backend & GET /report/csv returns valid CSV               & Pass \\
Backend & PATCH /sessions/\{id\}/end closes session       & Pass \\
Backend & POST /heartbeat on closed session (409)         & Pass \\
Backend & GET /dashboard returns HTML                     & Pass \\
\midrule
Router  & GET / returns check-in form HTML                & Pass \\
Router  & GET /ping returns pong                          & Pass \\
Router  & GET /session returns JSON or 204                & Pass \\
Router  & POST /checkin handles unknown roll number       & Pass \\
Router  & POST /heartbeat rejects wrong session\_id (409) & Pass \\
Router  & POST /heartbeat rejects empty body (400)        & Pass \\
Router  & GET /success has JS keep-alive script           & Pass \\
Router  & Captive portal probes return 204                & Pass \\
\midrule
\multicolumn{2}{l}{\textbf{Total}} & \textbf{20/20} \\
\bottomrule
\end{tabular}
\end{table}

\subsection{Load Test}

\texttt{load\_test.py} simulates 35 concurrent students: each thread
enrolls a unique token, sends 10 heartbeat rounds at 0.5~s intervals,
then the final report is fetched and validated.
Table~\ref{tab:load} shows the results.

\begin{table}[h]
\caption{Load test: 35 concurrent students, FastAPI + SQLite, loopback}
\label{tab:load}
\begin{tabular}{lr}
\toprule
Metric & Value \\
\midrule
p50 heartbeat latency      & $<15$~ms  \\
p95 heartbeat latency      & $<50$~ms  \\
p99 heartbeat latency      & $<120$~ms \\
Report generation (35 students) & $<200$~ms \\
Classification accuracy    & 100\% (35/35 PRESENT) \\
\bottomrule
\end{tabular}
\end{table}

Real LAN latency adds 1--5~ms per hop, still well within the 500~ms
target.  For deployments requiring higher concurrency, replacing SQLite
with PostgreSQL eliminates write serialisation with no API changes
(SQLAlchemy abstracts the driver).

\subsection{Physical Device Test}

A student phone was connected to the professor's Android hotspot (gateway
IP 10.116.213.247 in this run).  The ping sweep detected the device
within one 5~s scan cycle at a measured RTT of 277~ms.  After navigating
to the captive portal and submitting roll number CS001, the Router App
dashboard reflected one checked-in student.  JavaScript heartbeats
appeared in the server log at 30~s intervals.  The attendance report at
session end correctly classified the student as PRESENT.

The anti-proxy guards were also verified: re-submitting a different roll
number from the same device produced the Guard~1 rejection message;
submitting CS001 from a second device produced the Guard~2 rejection
message.

%% -----------------------------------------------------------
\section{Limitations and Future Work}
%% -----------------------------------------------------------

\paragraph{Hotspot client cap.}
The Android \texttt{LocalOnlyHotspot} API limits connected devices to
8--10 on most OEM firmware.  Larger classes require a dedicated travel
router (e.g., TP-Link TL-WR802N) or a Raspberry Pi in AP mode.

\paragraph{Residual proxy risk.}
Guard~1 and Guard~2 prevent the most common forms of proxy check-in,
but cannot stop a student from physically handing their phone to a
friend.  Eliminating this risk requires a biometric or one-time PIN at
check-in time, which would necessitate a thin native app on the
student's device.

\paragraph{Captive portal auto-detection.}
Android's captive-portal detection varies by OEM and Android version.
DNS interception (resolving all hostnames to the hotspot IP) would
force automatic redirect but requires root access or a custom DNS
server on the hotspot.

\paragraph{Clock skew.}
Presence-duration computation depends on consistent wall-clock time.
Devices without NTP may skew up to 2~s, negligible for 75-minute
sessions but potentially material for short quizzes.

%% -----------------------------------------------------------
\section{Conclusion}
%% -----------------------------------------------------------

This paper presented a WiFi-based classroom attendance system that runs
on a single commodity Android smartphone, requires no dedicated student
app, and stores no personally identifiable network information.  The
dual-layer presence signal --- ICMP ping sweep for network-layer
connectivity and JavaScript heartbeats for application-layer
confirmation --- provides robust presence-duration estimation without
any hardware beyond the professor's personal device.  Two server-side
guards prevent the most common forms of proxy check-in.  All 20
functional tests pass on a live device and a 35-student load test
confirms sub-15~ms median heartbeat latency.  The system is ready for
pilot deployment in a real classroom.

%% -----------------------------------------------------------
\begin{thebibliography}{9}

\bibitem{demir2014}
M.~Demir, N.~Bhatt, D.~Mendez, H.~Mineno, and A.~Leon-Garcia,
``Passive WiFi Localization and Activity Monitoring,''
in \emph{Proc.\ IEEE ICC}, 2014.

\bibitem{hwang2017}
J.-H.~Hwang and L.-C.~Wang,
``WiFi-based Classroom Attendance Tracking Using DHCP Logs,''
in \emph{Proc.\ IEEE COMPSAC}, 2017.

\bibitem{nanohttpd}
NanoHTTPD Contributors,
``NanoHTTPD --- a tiny web server in Java,''
GitHub, 2023. [Online]. Available: \url{https://github.com/NanoHttpd/nanohttpd}

\bibitem{fastapi}
S.~Ramirez,
``FastAPI,''
GitHub, 2023. [Online]. Available: \url{https://github.com/tiangolo/fastapi}

\bibitem{room}
Android Developers,
``Room Persistence Library,''
Android Documentation, 2023. [Online]. Available:
\url{https://developer.android.com/training/data-storage/room}

\bibitem{localhotspot}
Android Developers,
``WifiManager.LocalOnlyHotspotCallback,''
\emph{Android API Reference}, 2023.

\end{thebibliography}

\end{document}
